\documentclass[12pt,a4paper]{article}

\usepackage[T1]{fontenc}
\usepackage[utf8]{inputenc}
\usepackage[french]{babel}
\usepackage[margin=2.3cm, top=2.8cm, bottom=2.8cm]{geometry}
\usepackage{hyperref}
\usepackage{xcolor}
\usepackage{tcolorbox}
\usepackage{tikz}
\usepackage{booktabs}
\usepackage{longtable}
\usepackage{enumitem}
\usepackage{lmodern}
\usepackage{microtype}
\usepackage{titlesec}
\usepackage{fancyhdr}
\usepackage{float}
\usepackage{multirow}
\usepackage{array}
\usepackage{caption}

\usetikzlibrary{
  shapes.geometric, shapes.multipart,
  arrows.meta, positioning, fit,
  backgrounds, calc
}

%──────────────────────────────
% COULEURS
%──────────────────────────────
\definecolor{agentblue}{RGB}{25, 80, 150}
\definecolor{agentgreen}{RGB}{30, 130, 60}
\definecolor{agentred}{RGB}{170, 30, 30}
\definecolor{agentorange}{RGB}{200, 110, 0}
\definecolor{agentpurple}{RGB}{100, 40, 140}
\definecolor{agentgray}{RGB}{245, 245, 248}
\definecolor{agentlightblue}{RGB}{218, 232, 252}
\definecolor{agentlightgreen}{RGB}{218, 245, 225}
\definecolor{agentlightred}{RGB}{252, 225, 225}
\definecolor{agentlightyellow}{RGB}{255, 248, 218}
\definecolor{agentlightpurple}{RGB}{238, 225, 252}
\definecolor{blockred}{RGB}{210, 50, 50}

%──────────────────────────────
% BOÎTES
%──────────────────────────────
\tcbuselibrary{skins, breakable}

\newtcolorbox{forcebox}[1][]{
  colback=agentlightgreen, colframe=agentgreen,
  fonttitle=\bfseries, title=#1, breakable, arc=5pt,
  left=8pt, right=8pt, top=5pt, bottom=5pt
}

\newtcolorbox{weaknessbox}[1][]{
  colback=agentlightred, colframe=blockred,
  fonttitle=\bfseries, title=#1, breakable, arc=5pt,
  left=8pt, right=8pt, top=5pt, bottom=5pt
}

\newtcolorbox{blockerbox}[1][]{
  colback=agentlightred!60, colframe=blockred, boxrule=1.2pt,
  fonttitle=\bfseries, title=\faExclamation~#1, breakable, arc=5pt,
  left=8pt, right=8pt, top=5pt, bottom=5pt
}

\newtcolorbox{notebox}{
  colback=agentlightyellow, colframe=agentorange,
  fonttitle=\bfseries, breakable, arc=5pt,
  left=8pt, right=8pt, top=5pt, bottom=5pt
}

\newtcolorbox{diagbox}[1][]{
  colback=agentlightblue, colframe=agentblue,
  fonttitle=\bfseries, title=#1, breakable, arc=5pt,
  left=8pt, right=8pt, top=5pt, bottom=5pt
}

%──────────────────────────────
% EN-TÊTES
%──────────────────────────────
\pagestyle{fancy}
\fancyhf{}
\fancyhead[L]{\textcolor{agentblue}{\textbf{AgentActuariat}} \textcolor{gray}{— Analyse diagnostique}}
\fancyhead[R]{\textcolor{gray}{\small Avril 2026}}
\fancyfoot[C]{\thepage}
\renewcommand{\headrulewidth}{0.5pt}

\titleformat{\section}{\Large\bfseries\color{agentblue}}{}{0em}{\thesection\quad}[\vspace{-2pt}\rule{\linewidth}{0.6pt}\vspace{2pt}]
\titleformat{\subsection}{\large\bfseries\color{agentblue!80!black}}{}{0em}{\thesubsection\quad}
\titleformat{\subsubsection}{\normalsize\bfseries\color{agentblue!60!black}}{}{0em}{\thesubsubsection\quad}

\hypersetup{
  colorlinks=true, linkcolor=agentblue, urlcolor=agentblue,
  pdftitle={AgentActuariat — Analyse diagnostique},
  pdfauthor={AgentActuariat}
}

%══════════════════════════════════════════════════════════════════
\begin{document}
%══════════════════════════════════════════════════════════════════

\begin{titlepage}
\centering
\vspace*{1.5cm}
{\fontsize{30}{36}\bfseries\color{agentblue} AgentActuariat\par}
\vspace{0.4cm}
{\Large\color{gray} Analyse diagnostique de l'architecture\par}
\vspace{0.3cm}
{\large\color{agentred}\textit{Pourquoi le système peine à produire un rapport professionnel}\par}
\vspace{1.5cm}
\rule{0.7\linewidth}{0.6pt}
\vspace{1cm}

\begin{tcolorbox}[colback=agentgray, colframe=agentblue!25,
                  width=0.85\textwidth, arc=6pt]
\small
\textbf{Objet :} Ce document décrit en français le fonctionnement complet
d'AgentActuariat, met en avant ses forces réelles, et identifie précisément
les points qui bloquent la qualité du rapport actuariel produit.

\medskip
\textbf{Approche :} Diagnostique. Pour chaque composant, on indique ce qui
fonctionne, ce qui pose problème, et l'impact concret sur le livrable.

\medskip
\textbf{Public :} Marc, pour comprendre où agir en priorité.
\end{tcolorbox}

\vfill
\begin{center}
\begin{tabular}{ll}
\toprule
\textbf{Système}        & AgentActuariat \\
\textbf{Architecture}   & LangGraph multi-agents \\
\textbf{Modèle LLM}     & GPT-4o (OpenAI) \\
\textbf{Livrable cible} & Rapport PDF de table de mortalité d'expérience \\
\bottomrule
\end{tabular}
\end{center}

\vspace{1cm}
{\footnotesize\color{gray} Avril 2026 · Version 2.0\par}
\end{titlepage}

\tableofcontents
\newpage

%══════════════════════════════════════════════════════════════════
\section{Synthèse exécutive — Le diagnostic en une page}
%══════════════════════════════════════════════════════════════════

\begin{diagbox}[Ce que fait le système]
À partir d'un fichier CSV (portefeuille assurantiel) et d'un rapport de référence,
AgentActuariat doit produire un rapport actuariel professionnel décrivant la
construction d'une table de mortalité d'expérience. Il enchaîne pour cela :

\begin{enumerate}[leftmargin=2em, itemsep=2pt]
  \item Une compréhension de la demande (MasterAgent).
  \item Une chaîne de calculs actuariels (BuilderAgent et ses 10 outils).
  \item Une rédaction automatisée du rapport (Pipeline en 6 étapes).
  \item Une validation finale du rapport (LLM juge).
\end{enumerate}
\end{diagbox}

\begin{forcebox}[Les forces réelles du système]
\begin{itemize}[leftmargin=1.5em, itemsep=2pt]
  \item Le \textbf{moteur de calcul actuariel} est solide (10 outils couvrant exposition,
        lissage, validation, benchmarking, régressions Cox et logit).
  \item Les \textbf{tableaux et graphiques} sont rendus de façon déterministe, sans LLM —
        donc reproductibles et sans hallucination.
  \item Un \textbf{validateur de traçabilité} vérifie que chaque chiffre cité dans le texte
        existe bien dans les données calculées.
  \item La \textbf{mémoire persistante} permet de reprendre une étude entre sessions.
\end{itemize}
\end{forcebox}

\begin{weaknessbox}[Les blocages structurels qui dégradent le rapport]
\begin{enumerate}[leftmargin=1.5em, itemsep=3pt]
  \item Le rapport de référence n'est \textbf{jamais lu en entier} : seuls 600 caractères
        par section sont injectés via RAG. Le style et la structure du modèle ne sont
        donc pas vraiment transmis.
  \item Le LLM rédige avec un \textbf{plafond de 1\,200 tokens par section}, souvent
        insuffisant pour une analyse actuarielle dense.
  \item Les étapes de validation (sufficience des données, qualité finale) tournent à
        \textbf{température 1.0 par défaut}, ce qui rend les verdicts non reproductibles.
  \item Le WriterNode \textbf{ne réfléchit pas avant de rédiger} : il déclenche le pipeline
        sans aucune analyse contextuelle des spécificités du portefeuille.
  \item Chaque section est rédigée \textbf{de façon indépendante} : aucune cohérence narrative
        n'est imposée d'une section à l'autre.
  \item Le validateur de traçabilité vérifie \textbf{les chiffres uniquement}, pas le sens :
        le LLM peut produire une interprétation fausse d'un chiffre correct.
\end{enumerate}
\end{weaknessbox}

\newpage

%══════════════════════════════════════════════════════════════════
\section{Vue d'ensemble du fonctionnement}
%══════════════════════════════════════════════════════════════════

\subsection{Schéma global}

\begin{figure}[H]
\centering
\begin{tikzpicture}[
  node distance=0.8cm and 1.5cm,
  every node/.style={font=\small},
  data/.style={cylinder, shape border rotate=90, aspect=0.3,
               draw=gray, fill=agentgray, minimum width=2cm,
               minimum height=1cm, font=\small},
  agent/.style={rectangle, rounded corners=6pt, draw=agentblue,
                fill=agentlightblue, minimum width=3.2cm,
                minimum height=1cm, text centered, font=\small\bfseries},
  process/.style={rectangle, rounded corners=4pt, draw=agentpurple,
                  fill=agentlightpurple, minimum width=3.2cm,
                  minimum height=1cm, text centered, font=\small\bfseries},
  output/.style={rectangle, rounded corners=6pt, draw=agentred,
                 fill=agentlightred, minimum width=3.2cm,
                 minimum height=1cm, text centered, font=\small\bfseries},
  arr/.style={-{Stealth[length=6pt]}, line width=1pt, color=agentblue}
]

\node[data]                  (csv)     {CSV};
\node[agent, right=of csv]   (master)  {MasterAgent\\ \footnotesize Comprend, route};
\node[agent, right=of master](builder) {BuilderAgent\\ \footnotesize Calcule};
\node[process, below=of builder] (pipeline) {Pipeline rapport\\ \footnotesize 6 étapes};
\node[output, below=of pipeline] (pdf) {PDF rapport\\ final};
\node[data, below=of master] (memory) {Mémoire\\persistante};
\node[data, left=of pipeline, xshift=-0.5cm] (yaml) {Template\\YAML};
\node[data, right=of pipeline, xshift=0.5cm] (rag) {Rapport\\exemple\\(RAG)};

\draw[arr] (csv)     -- (master);
\draw[arr] (master)  -- (builder);
\draw[arr] (builder) -- (pipeline);
\draw[arr] (pipeline)-- (pdf);
\draw[arr] (yaml)    -- (pipeline);
\draw[arr] (rag)     -- (pipeline);
\draw[arr, dashed, gray] (memory.north) -- (master.south);
\draw[arr, dashed, gray] (builder.south) -- (memory.north east);

\end{tikzpicture}
\caption{Vue d'ensemble du flux de bout en bout}
\end{figure}

\subsection{Composants et leur rôle}

\begin{center}
\begin{tabular}{p{3.2cm} p{6.5cm} p{3cm}}
\toprule
\textbf{Composant} & \textbf{Rôle} & \textbf{Utilise un LLM} \\
\midrule
Interface Dash    & Dialogue avec l'utilisateur, upload du CSV     & Non \\
\addlinespace
MasterAgent       & Comprend l'intention, oriente la conversation  & Oui (3 à 4 appels) \\
\addlinespace
BuilderAgent      & Choisit et enchaîne les outils de calcul       & Oui (boucle ReAct) \\
\addlinespace
Outils actuariels & Calculs déterministes (exposition, lissage, etc.) & Non \\
\addlinespace
WriterNode        & Déclenche le pipeline de rédaction             & Non \\
\addlinespace
Pipeline rapport  & Rédige et assemble le PDF en 6 étapes          & Partiel \\
\addlinespace
Mémoire           & Conserve l'état entre les tours et les sessions & Non \\
\bottomrule
\end{tabular}
\end{center}

\newpage

%══════════════════════════════════════════════════════════════════
\section{Le MasterAgent — Le chef d'orchestre}
%══════════════════════════════════════════════════════════════════

\subsection{Ce qu'il fait}

Le MasterAgent est le seul composant autorisé à décider qui parle ensuite.
À chaque tour, il :

\begin{enumerate}[leftmargin=2em]
  \item Classifie l'intention de l'utilisateur (calcul, rapport, question).
  \item Au premier tour : désambiguïse les colonnes du CSV et le plan d'étude.
  \item Extrait le plan d'étude depuis l'historique conversationnel.
  \item Décide vers qui router : BuilderAgent, WriterNode, ou réponse directe.
  \item Répond conversationnellement aux questions simples.
\end{enumerate}

\subsection{Logique de décision}

\begin{figure}[H]
\centering
\begin{tikzpicture}[
  node distance=0.4cm,
  decision/.style={diamond, draw=agentorange, fill=agentlightyellow,
                   minimum width=4cm, aspect=2.5,
                   text centered, font=\footnotesize, inner sep=1pt},
  action/.style={rectangle, rounded corners=4pt, draw=agentblue,
                 fill=agentlightblue, minimum width=2.8cm,
                 minimum height=0.7cm, text centered, font=\footnotesize\bfseries},
  arr/.style={-{Stealth[length=4pt]}, line width=0.6pt},
  lbl/.style={font=\tiny\itshape, fill=white, inner sep=2pt}
]

\node[decision] (d1) {Le rapport est-il fini ?};
\node[action, right=2cm of d1] (e1) {Terminer};
\node[decision, below=0.5cm of d1] (d2) {Les calculs sont-ils finis et un rapport demandé ?};
\node[action, right=2cm of d2] (e2) {Aller au Writer};
\node[decision, below=0.5cm of d2] (d3) {Le Writer manque de données ?};
\node[action, right=2cm of d3] (e3) {Aller au Builder};
\node[decision, below=0.5cm of d3] (d4) {Quelle intention ?};
\node[action, right=2.5cm of d4, yshift=0.6cm] (e4a) {Builder};
\node[action, right=2.5cm of d4] (e4b) {Writer};
\node[action, right=2.5cm of d4, yshift=-0.6cm] (e4c) {Réponse directe};

\draw[arr] (d1) -- node[lbl, above]{Oui} (e1);
\draw[arr] (d1) -- node[lbl, right]{Non} (d2);
\draw[arr] (d2) -- node[lbl, above]{Oui} (e2);
\draw[arr] (d2) -- node[lbl, right]{Non} (d3);
\draw[arr] (d3) -- node[lbl, above]{Oui} (e3);
\draw[arr] (d3) -- node[lbl, right]{Non} (d4);
\draw[arr] (d4.east) -- node[lbl, above]{calcul} (e4a.west);
\draw[arr] (d4.east) -- node[lbl, above]{rédaction} (e4b.west);
\draw[arr] (d4.east) -- node[lbl, below]{question} (e4c.west);

\end{tikzpicture}
\caption{Arbre de décision du MasterAgent à chaque tour}
\end{figure}

\subsection{Forces}

\begin{forcebox}[Ce qui fonctionne bien]
\begin{itemize}[leftmargin=1.5em]
  \item La séparation claire des rôles : un seul composant décide du routing,
        ce qui évite les conflits de décision.
  \item La désambiguïsation au premier tour évite que le système se lance sur
        des mauvaises colonnes.
  \item Le system prompt est enrichi dynamiquement avec l'état du calcul
        (résultats déjà disponibles, SMR global, mapping des colonnes).
\end{itemize}
\end{forcebox}

\subsection{Faiblesses}

\begin{weaknessbox}[Points qui posent problème]
\begin{itemize}[leftmargin=1.5em]
  \item Le contenu exact des consignes données au LLM (system prompt) n'a pas été
        audité — c'est un angle mort. On ne sait pas vraiment ce que le MasterAgent
        a appris à faire.
  \item La classification d'intention se fait avec seulement 200 tokens en sortie,
        ce qui force le LLM à être très concis et peut produire des classifications
        approximatives sur des demandes nuancées.
  \item Aucun mécanisme de confirmation avant de lancer un calcul lourd : si la
        désambiguïsation se trompe, l'utilisateur doit interrompre.
\end{itemize}
\end{weaknessbox}

\newpage

%══════════════════════════════════════════════════════════════════
\section{Le BuilderAgent — Le calculateur}
%══════════════════════════════════════════════════════════════════

\subsection{Ce qu'il fait}

Le BuilderAgent est un \textbf{agent ReAct} : il alterne raisonnement et action.
Il dispose de 10 outils actuariels qu'il peut appeler dans l'ordre qu'il choisit.

\begin{figure}[H]
\centering
\begin{tikzpicture}[
  node distance=1.2cm,
  box/.style={rectangle, rounded corners=5pt, draw=agentgreen,
              fill=agentlightgreen, minimum width=4cm,
              minimum height=1cm, text centered, font=\small\bfseries},
  arr/.style={-{Stealth[length=5pt]}, line width=0.8pt, color=agentgreen},
  lbl/.style={font=\footnotesize\itshape, fill=white, inner sep=2pt}
]

\node[box] (llm) {LLM GPT-4o\\ \footnotesize raisonne et choisit};
\node[box, below=of llm] (tool) {Outil actuariel\\ \footnotesize calcule};
\node[box, right=2.5cm of llm] (data) {Mémoire\\ \footnotesize résultats};

\draw[arr, bend left=20] (llm) to node[lbl, right]{appel d'outil} (tool);
\draw[arr, bend left=20] (tool) to node[lbl, left]{résultat} (llm);
\draw[arr] (tool) -- node[lbl, below]{stocke} (data.south);

\node[font=\footnotesize\itshape, color=agentred,
      below=0.3cm of tool, text width=8cm, align=center]
  {Boucle limitée à 5 itérations pour éviter les boucles infinies};

\end{tikzpicture}
\caption{Boucle de raisonnement-action du BuilderAgent}
\end{figure}

\subsection{Les 10 outils disponibles}

\begin{center}
\begin{tabular}{p{4.5cm} p{4cm} p{4.5cm}}
\toprule
\textbf{Outil} & \textbf{Donnée d'entrée} & \textbf{Résultat produit} \\
\midrule
Calcul de l'exposition  & DataFrame du portefeuille & Années-personnes par âge \\
Taux bruts ($q_x$)      & Table d'exposition        & Probabilités brutes par âge \\
Lissage                 & Taux bruts                & Probabilités lissées \\
Diagnostics             & Table d'exposition        & Crédibilité, recommandations \\
Validation statistique  & Exposition et lissage     & Intervalles de confiance, $\chi^2$ \\
Benchmarking            & Exposition et lissage     & SMR global, abattement \\
Régression de Cox       & DataFrame et exposition   & Hazard ratio par sexe \\
Régression logit        & Lissage et table de référence & Pente, $R^2$ \\
Comparaison antérieure  & Lissage                   & Dérive vs étude précédente \\
Synthèse de portefeuille& DataFrame                 & Statistiques descriptives \\
\bottomrule
\end{tabular}
\end{center}

\subsection{Forces}

\begin{forcebox}[Ce qui fonctionne bien]
\begin{itemize}[leftmargin=1.5em]
  \item \textbf{Tous les calculs sont déterministes} : aucune hallucination possible
        sur les chiffres produits par les outils.
  \item Le LLM \textbf{adapte le system prompt à la phase} : court après une exécution
        d'outil, complet quand il faut planifier la suite.
  \item La boucle est \textbf{bornée à 5 itérations}, ce qui empêche les boucles
        infinies coûteuses.
  \item Les résultats sont \textbf{accumulés dans une mémoire partagée} (\texttt{data\_store})
        accessible aux autres composants.
\end{itemize}
\end{forcebox}

\subsection{Faiblesses}

\begin{weaknessbox}[Points qui posent problème]
\begin{itemize}[leftmargin=1.5em]
  \item \textbf{Pas de garde-fou proactif sur l'ordre des outils.} Le LLM peut
        appeler le benchmarking avant le lissage : l'outil retournera une erreur,
        mais cela consomme une itération sur les 5 disponibles.
  \item Le seul point bloquant identifié dans les outils est le contrôle de
        monotonie du lissage — tout le reste passe en mode "passif".
  \item L'historique est limité à 20 messages : au-delà, un résumé automatique
        prend le relais, ce qui peut faire perdre des nuances importantes du dialogue.
  \item La température du LLM est laissée à la valeur par défaut (1.0), ce qui
        peut produire des choix d'outils différents pour la même demande.
\end{itemize}
\end{weaknessbox}

\newpage

%══════════════════════════════════════════════════════════════════
\section{Le WriterNode et le Pipeline de rapport}
%══════════════════════════════════════════════════════════════════

\subsection{Le WriterNode — Un simple déclencheur}

Le WriterNode porte un nom trompeur : \textbf{il ne rédige rien lui-même}.
Il appelle directement le pipeline de rapport et retourne le chemin du PDF produit.

\begin{blockerbox}[Premier blocage majeur identifié]
Le WriterNode \textbf{n'effectue aucun appel LLM avant de déclencher le pipeline}.
Il n'y a donc aucune analyse intelligente du contexte avant la rédaction :
\begin{itemize}[leftmargin=1.5em, itemsep=0pt]
  \item Pas d'identification des particularités du portefeuille.
  \item Pas de détection des âges à faible crédibilité avant rédaction.
  \item Pas d'adaptation du ton ou de la structure selon les résultats.
  \item Pas de briefing transmis aux étapes suivantes du pipeline.
\end{itemize}
Conséquence : le pipeline traite tous les portefeuilles de la même manière,
indépendamment de leurs spécificités.
\end{blockerbox}

\subsection{Le pipeline en 6 étapes}

\begin{figure}[H]
\centering
\begin{tikzpicture}[
  node distance=0.5cm,
  step/.style={rectangle, rounded corners=4pt, draw=agentblue,
               fill=agentlightblue, minimum width=10cm,
               minimum height=0.85cm, text centered, font=\small},
  llmstep/.style={step, draw=agentpurple, fill=agentlightpurple},
  ragstep/.style={step, draw=agentgreen, fill=agentlightgreen},
  outstep/.style={step, draw=agentred, fill=agentlightred,
                  minimum width=5cm},
  arr/.style={-{Stealth[length=4pt]}, line width=0.7pt, color=agentblue}
]

\node[step]                (s1) {\textbf{1.} Préparation du plan (déterministe)};
\node[llmstep, below=of s1](s2) {\textbf{2.} Validation des données par LLM};
\node[ragstep, below=of s2](s3) {\textbf{3.} Recherche d'exemples par RAG};
\node[llmstep, below=of s3](s4) {\textbf{4.} Rédaction des sections par LLM};
\node[step,    below=of s4](s5) {\textbf{5.} Assemblage du PDF (déterministe)};
\node[llmstep, below=of s5](s6) {\textbf{6.} Validation finale par LLM};
\node[outstep, below=0.7cm of s6] (pdf) {\textbf{PDF livré}};

\foreach \a/\b in {s1/s2, s2/s3, s3/s4, s4/s5, s5/s6, s6/pdf}
  \draw[arr] (\a) -- (\b);

\end{tikzpicture}
\caption{Les 6 étapes du pipeline de rédaction}
\end{figure}

\subsubsection{Étape 1 — Préparation du plan}

Le système charge un template YAML qui décrit la structure du rapport
(sections, sous-sections, tableaux et graphiques attendus). Il y substitue
automatiquement les valeurs calculées (par exemple "exposition totale = 12 345 années-personnes").

\textbf{Force :} la structure du rapport est garantie et reproductible.\\
\textbf{Faiblesse :} la structure est figée, elle ne s'adapte pas au cas particulier.

\subsubsection{Étape 2 — Validation des données}

Un appel LLM vérifie que chaque section dispose des données nécessaires.
Si une donnée manque, le pipeline s'arrête et renvoie au BuilderAgent.

\begin{blockerbox}[Deuxième blocage : verdict non déterministe]
Cette étape utilise GPT-4o avec une \textbf{température par défaut de 1.0}.
Conséquence : le verdict (données suffisantes ou non) peut varier entre deux exécutions
identiques. Pour une étape de validation, c'est inacceptable.
\end{blockerbox}

\subsubsection{Étape 3 — Recherche d'exemples (RAG)}

Le système interroge une base vectorielle (ChromaDB) qui contient le rapport
de référence. Pour chaque section, il extrait des passages similaires destinés
à inspirer le style.

\begin{blockerbox}[Troisième blocage : RAG insuffisant]
Les passages récupérés sont \textbf{tronqués à 600 caractères}, soit l'équivalent
d'environ 100 mots. C'est très insuffisant pour transmettre le style et la
structure d'une section actuarielle qui fait typiquement 600 à 900 mots.

De plus, \textbf{le rapport de référence n'est jamais lu en entier} par le système.
Le RAG ne récupère que des fragments isolés, sans cohérence narrative globale.
C'est probablement \textbf{le blocage le plus important} pour reproduire un style
professionnel.
\end{blockerbox}

\subsubsection{Étape 4 — Rédaction des sections}

C'est le cœur de la rédaction. Pour chaque section, en parallèle :

\begin{enumerate}[leftmargin=2em]
  \item Les tableaux et graphiques sont rendus de façon déterministe
        (ReportLab et matplotlib), avec injection des bonnes données.
  \item Le LLM rédige le texte narratif en s'appuyant sur le prompt préparé
        à l'étape 1, enrichi des exemples de l'étape 3 et des tableaux rendus.
  \item Le validateur de traçabilité vérifie que tous les chiffres cités sont
        bien issus des calculs (voir section dédiée).
\end{enumerate}

\textbf{Paramètres LLM :} température 0.4, plafond de 1\,200 tokens par section.

\begin{blockerbox}[Quatrième blocage : sections trop courtes]
Le plafond de \textbf{1\,200 tokens} (environ 850 mots) est suffisant pour les
sections courtes (préambule, conclusion) mais \textbf{trop juste pour les sections
denses} comme la construction de la table ou l'analyse des résultats.

Une section actuarielle professionnelle de qualité fait typiquement 800 à
1\,200 mots, soit 1\,500 à 2\,000 tokens. Le risque est que le LLM tronque ses
explications ou omette des éléments d'analyse.
\end{blockerbox}

\begin{blockerbox}[Cinquième blocage : pas de cohérence inter-sections]
Les 5 sections sont rédigées \textbf{en parallèle, indépendamment}.
Aucune section ne sait ce que les autres ont écrit. Conséquence :
\begin{itemize}[leftmargin=1.5em, itemsep=0pt]
  \item Risque de répétitions (même chiffre commenté deux fois différemment).
  \item Risque d'incohérences (interprétations divergentes du même résultat).
  \item Pas de fil narratif global.
\end{itemize}
\end{blockerbox}

\subsubsection{Étape 5 — Assemblage du PDF}

Une étape purement déterministe qui assemble les sections, les tableaux et
les graphiques dans l'ordre fixe défini par le template. Aucune surprise ici.

\subsubsection{Étape 6 — Validation finale}

Un appel LLM joue le rôle de "réviseur qualité actuariel senior". Il reçoit
un résumé de chaque section (limité à 300 mots) et émet un verdict :
\textit{ok}, \textit{minor} (anomalies mineures) ou \textit{major} (anomalies majeures).

Si le verdict est \textit{minor}, les sections concernées sont régénérées
une fois.

\begin{blockerbox}[Sixième blocage : validation finale aveugle]
Cette étape souffre de \textbf{deux limitations majeures} :

\begin{enumerate}[leftmargin=1.5em, itemsep=0pt]
  \item Comme à l'étape 2, la \textbf{température est laissée à 1.0 par défaut}.
        Le verdict est non reproductible.
  \item Le réviseur ne voit qu'\textbf{un résumé de 300 mots par section}, pas
        le texte intégral. Il ne peut donc pas détecter une mauvaise tournure de
        phrase, une interprétation erronée, ou une formulation non professionnelle.
\end{enumerate}
\end{blockerbox}

\newpage

%══════════════════════════════════════════════════════════════════
\section{Le validateur de traçabilité}
%══════════════════════════════════════════════════════════════════

\subsection{Ce qu'il fait}

Pour chaque section rédigée, le validateur extrait par expression régulière
tous les nombres présents dans le texte, puis vérifie qu'ils correspondent
à des valeurs calculées dans la mémoire (avec tolérance de 2 \%).

Si des nombres ne sont pas traçables, un retour ciblé est donné au LLM,
qui régénère la section une fois.

\subsection{Force majeure}

\begin{forcebox}[Une garantie forte sur les chiffres]
C'est l'une des composantes les plus solides du système. Aucun nombre inventé
ne devrait apparaître dans le rapport final. Les conversions usuelles
(fraction vers pourcent, pourcent vers millième) sont gérées, et une liste
blanche couvre les constantes statistiques universelles (1{,}96, 0{,}05, etc.).
\end{forcebox}

\subsection{Limite à connaître}

\begin{weaknessbox}[Ce que le validateur ne fait pas]
Le validateur \textbf{ne vérifie que les chiffres, pas le sens du texte}.
Le LLM peut parfaitement écrire :

\medskip
\textit{« Le SMR global de 0,87 indique que la mortalité du portefeuille
est supérieure à la référence. »}
\medskip

Le chiffre 0{,}87 est traçable, donc le validateur le valide. Mais
l'interprétation est \textbf{fausse} (un SMR inférieur à 1 indique une mortalité
\emph{inférieure}). \textbf{Aucun mécanisme du système ne détecte ce type d'erreur}.
\end{weaknessbox}

\newpage

%══════════════════════════════════════════════════════════════════
\section{Le système de mémoire}
%══════════════════════════════════════════════════════════════════

\subsection{Quatre couches superposées}

\begin{center}
\begin{tabular}{p{3.5cm} p{4.5cm} p{3cm} p{2.5cm}}
\toprule
\textbf{Couche} & \textbf{Contenu} & \textbf{Persistance} & \textbf{Survit au redémarrage} \\
\midrule
Mémoire de travail & État LangGraph complet & RAM & Non \\
Mémoire métier & Plan d'étude, mapping colonnes, résultats actuariels & JSON sur disque & Oui \\
Mémoire des données & DataFrame du portefeuille & Parquet sur disque & Oui \\
Mémoire conversationnelle & Historique des messages compactés & RAM & Non \\
\bottomrule
\end{tabular}
\end{center}

\subsection{Forces}

\begin{forcebox}[Ce qui fonctionne bien]
\begin{itemize}[leftmargin=1.5em]
  \item Les résultats actuariels coûteux à recalculer sont \textbf{persistés sur disque}.
  \item Le DataFrame brut n'est écrit qu'une seule fois, pas à chaque tour.
  \item Au-delà de 15 messages, un résumé automatique structuré prend le relais
        pour éviter de saturer le contexte du LLM.
\end{itemize}
\end{forcebox}

\subsection{Faiblesses}

\begin{weaknessbox}[Limites importantes à comprendre]
\begin{itemize}[leftmargin=1.5em]
  \item Seules les clés présentes dans une \textbf{liste blanche fermée} sont sauvegardées
        sur disque. Tout résultat intermédiaire non listé est perdu au prochain tour.
  \item Si le processus est redémarré au milieu d'une conversation, l'\textbf{historique
        des messages est perdu}. Seuls les résultats actuariels sont récupérables.
  \item Le résumé automatique d'historique peut perdre des nuances importantes du dialogue
        (par exemple une demande spécifique de l'utilisateur sur le ton du rapport).
\end{itemize}
\end{weaknessbox}

\newpage

%══════════════════════════════════════════════════════════════════
\section{Diagnostic — Pourquoi le rapport n'est pas à la hauteur}
%══════════════════════════════════════════════════════════════════

\subsection{Hiérarchie des causes}

L'analyse fait ressortir une hiérarchie claire des problèmes, du plus impactant
au moins critique :

\begin{figure}[H]
\centering
\begin{tikzpicture}[
  node distance=0.4cm,
  block/.style={rectangle, rounded corners=5pt,
                minimum width=12cm, minimum height=1.1cm,
                text centered, font=\small\bfseries, draw, line width=1pt},
  critique/.style={block, fill=blockred!20, draw=blockred},
  haut/.style={block, fill=agentorange!20, draw=agentorange},
  moyen/.style={block, fill=agentlightyellow, draw=agentorange!60}
]

\node[critique] (b1) {1. Rapport de référence non lu en entier (RAG limité à 600 caractères)};
\node[critique, below=of b1] (b2) {2. WriterNode sans intelligence contextuelle};
\node[critique, below=of b2] (b3) {3. Pas de cohérence inter-sections (rédaction parallèle)};
\node[haut,     below=of b3] (b4) {4. Validateur final aveugle (résumé de 300 mots, température 1.0)};
\node[haut,     below=of b4] (b5) {5. Plafond de 1\,200 tokens trop bas pour sections denses};
\node[haut,     below=of b5] (b6) {6. Validation des données non déterministe (température 1.0)};
\node[moyen,    below=of b6] (b7) {7. Validateur de traçabilité ne vérifie que les chiffres};
\node[moyen,    below=of b7] (b8) {8. Garde-fous passifs sur l'ordre des outils};

\end{tikzpicture}
\caption{Hiérarchie des blocages identifiés (rouge : critique, orange : élevé, jaune : moyen)}
\end{figure}

\subsection{Le blocage principal}

\begin{blockerbox}[Le système ne sait pas vraiment imiter le rapport de référence]
\textbf{C'est le diagnostic central.} Le système a été pensé pour produire un rapport
à partir de données calculées et d'un template structurel, mais le rapport de
référence (qui est censé être le modèle de style et de structure narrative)
n'est utilisé que de façon très superficielle :

\begin{itemize}[leftmargin=1.5em]
  \item Il est découpé en chunks indexés dans une base vectorielle.
  \item Pour chaque section à rédiger, le système récupère des passages
        similaires (limités à 600 caractères chacun).
  \item Aucune compréhension globale du rapport de référence n'est construite.
  \item Aucune extraction de la voix, du ton, des tournures récurrentes.
\end{itemize}

Conséquence : le rapport produit a la \textbf{structure du template YAML}, les
\textbf{chiffres des calculs actuariels}, mais \textbf{pas le style du modèle}.
C'est pourquoi le résultat semble "synthétique" plutôt que "professionnel".
\end{blockerbox}

\subsection{Les leviers d'action — Par ordre d'efficacité}

\begin{diagbox}[Ce qui apporterait le plus de qualité, dans l'ordre]

\subsubsection*{1. Construire une vraie compréhension du rapport de référence}

Plutôt qu'un RAG par fragments, faire une passe initiale unique où un LLM
lit le rapport de référence en entier et en extrait :
\begin{itemize}[leftmargin=1.5em, itemsep=0pt]
  \item La voix narrative (ton, registre, niveau de formalisme).
  \item Les tournures récurrentes par type de section.
  \item Les transitions habituelles entre sections.
  \item Les conventions de présentation des chiffres.
\end{itemize}
Ce "guide de style" serait injecté dans tous les prompts de rédaction.

\subsubsection*{2. Ajouter une étape de briefing contextuel avant la rédaction}

Avant de lancer le pipeline, un appel LLM léger analyse le \texttt{data\_store}
et produit une note de cadrage qui identifie :
\begin{itemize}[leftmargin=1.5em, itemsep=0pt]
  \item Les particularités du portefeuille à mettre en avant.
  \item Les zones de faiblesse statistique à signaler.
  \item Le message principal que le rapport doit porter.
\end{itemize}

\subsubsection*{3. Introduire une cohérence inter-sections}

Plutôt qu'une rédaction parallèle, rédiger les sections en \textbf{séquence}, en
passant à chaque nouvelle section un résumé de ce qui a déjà été écrit et de
quoi parlera la section suivante.

\subsubsection*{4. Augmenter le plafond de tokens et fixer les températures}

Corrections rapides :
\begin{itemize}[leftmargin=1.5em, itemsep=0pt]
  \item Plafond à 2\,000 tokens pour les sections \textit{construction} et \textit{analysis}.
  \item Température à 0,1 sur les étapes 2 et 6 (validation).
\end{itemize}

\subsubsection*{5. Validateur final voyant le texte complet}

Au lieu d'un résumé de 300 mots, donner au validateur final l'intégralité du
texte, avec une grille d'évaluation explicite (cohérence, professionnalisme,
exactitude des interprétations).

\end{diagbox}

\newpage

%══════════════════════════════════════════════════════════════════
\section{Synthèse — Ce qu'il faut retenir}
%══════════════════════════════════════════════════════════════════

\begin{center}
\begin{tabular}{p{7cm} p{7cm}}
\toprule
\textbf{\color{agentgreen} Forces réelles} & \textbf{\color{blockred} Blocages structurels} \\
\midrule
Calcul actuariel complet et solide
& Le rapport de référence n'est pas vraiment lu \\
\addlinespace
Tableaux et graphiques déterministes
& Pas d'intelligence contextuelle avant rédaction \\
\addlinespace
Validateur de traçabilité des chiffres
& Sections rédigées indépendamment, sans cohérence \\
\addlinespace
Mémoire persistante des résultats
& Plafond de tokens trop bas pour sections denses \\
\addlinespace
Boucle ReAct bornée à 5 itérations
& Validations à température 1.0 (non déterministes) \\
\addlinespace
Désambiguïsation au premier tour
& Validateur final voit seulement un résumé \\
\addlinespace
Mode pas-à-pas avec approbation
& Validation des chiffres mais pas du sens \\
\bottomrule
\end{tabular}
\end{center}

\bigskip

\begin{diagbox}[Conclusion]
Le système est techniquement bien construit sur sa partie \textbf{calcul actuariel}.
Sa partie \textbf{rédaction de rapport} souffre en revanche d'un défaut de conception
majeur : elle traite la génération de texte comme une suite de tâches indépendantes
plutôt que comme l'écriture d'un document cohérent inspiré d'un modèle.

\medskip
Pour produire un rapport vraiment professionnel, il faut faire passer le système
d'une logique de \textbf{remplissage de template} à une logique de \textbf{rédaction
contextuelle inspirée d'un exemple}. C'est un changement d'architecture, pas
seulement un réglage de paramètres.
\end{diagbox}

\vfill
\begin{center}
\rule{0.5\linewidth}{0.4pt}\\[4pt]
{\footnotesize\color{gray}
Document produit à partir de \texttt{docs/AGENT\_ARCHITECTURE.md}\\
AgentActuariat — Analyse diagnostique — Avril 2026\par}
\end{center}

\end{document}
